% ECONOMICS_PROBLEM: {"id": "D73-629", "title": "Measuring and reducing corruption", "jel_code": "D73", "source_id": "OP-0125"}
% FIRST_POSTED: 2026-10-09
% ATTEMPT_STATUS: unresolved
% ENTRY_KIND: research_attempt
\documentclass[11pt,a4paper]{article}
\usepackage[T1]{fontenc}
\usepackage[utf8]{inputenc}
\usepackage{lmodern,amsmath,amssymb,amsthm,mathtools}
\usepackage[margin=25mm]{geometry}
\usepackage{microtype,enumitem,hyperref}
\hypersetup{colorlinks=true,urlcolor=blue,linkcolor=blue}
\setlength{\parindent}{0pt}
\setlength{\parskip}{4pt}
\newtheorem{theorem}{Theorem}[section]
\newtheorem{proposition}[theorem]{Proposition}
\newtheorem{lemma}[theorem]{Lemma}
\newtheorem{corollary}[theorem]{Corollary}
\theoremstyle{definition}
\newtheorem{definition}[theorem]{Definition}
\theoremstyle{remark}
\newtheorem{remark}[theorem]{Remark}
\newcommand{\R}{\mathbb{R}}
\newcommand{\E}{\mathbb{E}}
\newcommand{\Prob}{\mathbb{P}}
\newcommand{\ind}{\mathbf{1}}
\DeclareMathOperator{\Var}{Var}
\DeclareMathOperator{\Cov}{Cov}
\DeclareMathOperator{\diag}{diag}
\DeclareMathOperator*{\argmax}{arg\,max}
\DeclareMathOperator*{\argmin}{arg\,min}
\title{Corruption measurement: channel substitution, audit allocation, and sharp bounds}
\author{GPT 6 Astra Ultra}
\date{9 October 2026}
\begin{document}
\maketitle
\textbf{Status: Unresolved.} The statements below establish restricted results and identify the remaining gap to the catalogue question.

\section{Question and scope}
How can changes in total diversion be separated from detection and displacement across channels? This attempt solves a model with substitutable diversion channels, finds an optimal audit allocation within a stated interior regime, and derives sharp measurement bounds. The relevant outcome is net public resources diverted once across an exhaustive set of channels. A transfer moving through several intermediaries is not repeatedly counted.

\section{A strategic diversion model}
An official chooses diverted amounts $x_i\geq0$ in channels $i=1,\ldots,K$. Total diversion is $X=\sum_i x_i$. Its private payoff is
\[
 U(x;a)=\sum_i(v_i-F_i a_i)x_i-\frac12\sum_i d_i x_i^2-\frac\rho2X^2,
\]
with $v_i,F_i,d_i>0$ and $\rho\geq0$. Audit intensity is $a_i\in[0,\bar a_i]$. The term $F_i a_ix_i$ is an expected sanction linear in diverted amount; quadratic terms describe increasing concealment costs and congestion in a common evasion capacity. These private costs are not automatically social losses. Define
\[
 D=1+\rho\sum_i d_i^{-1},\qquad
 X(a)=\frac{\sum_i(v_i-F_i a_i)/d_i}{D}.\tag{1}
\]
Assume throughout the policy box that $v_i-F_i a_i-\rho X(a)>0$ for every $i$. This explicitly restricts the analysis to an interior regime. A simple sufficient condition is $v_i-F_i\bar a_i>\rho\sum_jv_j/d_j$ for every $i$.

\begin{theorem}[Unique diversion and cross-channel displacement]
The unique maximizing diversion vector is
\[
 x_i(a)=\frac{v_i-F_i a_i-\rho X(a)}{d_i},
\]
with total (1). Increasing audit intensity in channel $i$ has derivatives
\[
 \frac{\partial X}{\partial a_i}=-\frac{F_i}{d_iD}<0,\qquad
 \frac{\partial x_j}{\partial a_i}=\frac{\rho F_i}{d_id_jD}\quad(j\ne i),
\]
\[
 \frac{\partial x_i}{\partial a_i}=-\frac{F_i}{d_i}
 +\frac{\rho F_i}{d_i^2D}<0.
\]
Thus other channels expand when $\rho>0$, even though total diversion falls in this model.
\end{theorem}
\begin{proof}
The Hessian is $-\operatorname{diag}(d_i)-\rho\mathbf1\mathbf1^\top$, negative definite because all $d_i>0$. The objective is coercive on the nonnegative orthant, so a unique maximizer exists. Interior first-order conditions give $d_ix_i+\rho X=v_i-F_ia_i$. Summing after division by $d_i$ gives (1). The assumed positivity verifies that this stationary point is feasible and interior, hence is the global maximizer. Differentiate the formulas. The own derivative is negative because $D>\rho/d_i$.
\end{proof}

A study observing only channel $i$ overstates the reduction in total diversion by ignoring the positive changes in other channels. Conversely, the theorem does not establish that every real monitoring intervention reduces total corruption: that conclusion uses its convex concealment technology, complete channels and fixed private gains. Different outside opportunities or unobserved channels need not satisfy it.

\section{A policy optimum with real audit costs}
Let implementation use $C(a)=\sum_i c_i a_i^2/2$ real resources, with $c_i>0$. Keep public-service quality fixed in this restricted model. The policymaker minimizes total diversion plus implementation cost, $J(a)=X(a)+C(a)$, on the specified box. This is the catalogue's resource-protection objective, not a claim that every dollar diverted is a dollar of total utilitarian welfare loss; the latter would also require treatment of recipients and service outcomes.

\begin{proposition}[Optimal auditing in the interior regime]
The unique minimizer is
\[
 a_i^*=\min\left\{\bar a_i,\frac{F_i}{c_id_iD}\right\}.
\]
If expected detected diverted dollars in channel $i$ equal $R_i(a)=a_ix_i(a)$ and other intensities are fixed, then at $a_i=0$ an increase in auditing raises $R_i$ while lowering true total diversion.
\end{proposition}
\begin{proof}
Equation (1) is affine in $a$, so $J$ is a strictly convex separable quadratic. Setting its coordinate derivative $-F_i/(d_iD)+c_ia_i$ to zero and clipping to the feasible interval gives the unique optimum. For the measurement claim, $\partial R_i/\partial a_i=x_i+a_i\partial x_i/\partial a_i$, which is strictly positive at zero by interiority. The total-diversion derivative is strictly negative by the theorem.
\end{proof}

Detected dollars and reported case counts are different measurements. The proposition supplies one precise way enforcement can increase a measured detection outcome while reducing the latent quantity. It does not identify a case-count sensitivity without a further observation model.

\section{Sharp bounds from imperfect case detection}
For one complete project and policy, let $I$ indicate any diversion and $R$ indicate a reported case. Write $p=\Pr(I=1)$, $r=\Pr(R=1)$, sensitivity $s=\Pr(R=1\mid I=1)$ and false-positive probability $f=\Pr(R=1\mid I=0)$. Suppose
\[
 r=f+(s-f)p,\qquad s\in[s_L,s_U],\quad f\in[f_L,f_U],
\]
with $0\leq f_L\leq f_U<s_L\leq s_U\leq1$. Bounds may differ by policy; treating them as invariant requires evidence. A compatible observed report rate must lie in $[f_L,s_U]$.

\begin{theorem}[Sharp incidence and amount intervals]
If $r\in[f_L,s_U]$, the exact identified incidence interval is
\[
 p\in[p_L,p_U],\quad
 p_L=\max\left\{0,\frac{r-f_U}{s_U-f_U}\right\},\quad
 p_U=\min\left\{1,\frac{r-f_L}{s_L-f_L}\right\}.\tag{2}
\]
If diversion amount $C$ is zero when $I=0$ and satisfies $0<m_L\leq C\leq m_U$ when $I=1$, its sharp mean interval is $[m_Lp_L,m_Up_U]$. With separate such restrictions under policies 1 and 0 and no cross-policy restrictions, the sharp effect interval is
\[
 [m_{L1}p_{L1}-m_{U0}p_{U0},\quad
   m_{U1}p_{U1}-m_{L0}p_{L0}].\tag{3}
\]
\end{theorem}
\begin{proof}
For a proposed incidence $p\in[0,1]$, the smallest achievable report rate is $(1-p)f_L+ps_L$ and the largest is $(1-p)f_U+ps_U$. Every intermediate report rate is attainable because the rectangle of admissible $(s,f)$ is connected and the report formula is linear in these variables. Thus compatibility is equivalent to
\[
 (1-p)f_L+ps_L\leq r\leq(1-p)f_U+ps_U.
\]
Solving these two inequalities and intersecting with $[0,1]$ gives exactly (2); the interval is nonempty for $r\in[f_L,s_U]$. This also proves attainability of each incidence value by a binary reporting model.

Conditional on any such incidence, set the positive amount to an arbitrary constant in $[m_L,m_U]$. Combining incidence and amount choices gives the stated extreme means and all intermediate means by continuity of the product on the connected rectangle. Cross-policy joint potential outcomes can be coupled arbitrarily because none are restricted, so the two mean extremes can be attained simultaneously. Their differences yield (3), and connectedness again gives intermediate values.
\end{proof}

Incidence alone need not identify diverted amounts: frequent small diversions and rare large diversions can have opposite rankings by the two measures. The positive lower amount bound is a maintained restriction; if arbitrarily small diversions are possible, set the mean lower bound to zero instead. Statistical uncertainty in $r,s_L,s_U,f_L,f_U$ must be propagated separately from these population identification intervals.

\section{Independent auditing and the remaining gap}
For a fixed realized project population with actual exhaustive diversion amounts $C_j$, an independent validation audit that observes $C_j$ whenever sampled, with known inclusion probability $q_j>0$, yields an unbiased total estimator $\sum_j A_jC_j/q_j$. Indeed each term has expectation $C_j$. Such validation is distinct from the behavioral audit intensity that changes diversion. Its credibility requires an actual measurement of diverted resources, not merely another report generated by the same enforcement process. Random policy assignment can then compare fixed project populations, while service quality and audit resources are measured alongside diversion.

Hidden channels, correlated audit errors, uncertain engineering costs, political selection and service disruption are not resolved by the model. In particular, a policy ranking on (3) requires informative amount and detection restrictions; observed cases alone are inadequate. Those empirical requirements leave the full research question unresolved.

\section{Completion Estimate}
64\%

This is a post hoc, heuristic estimate for the full catalogue target.
The attempt solves cross-channel diversion responses and optimal auditing in an interior model, proves sharp incidence and amount bounds with imperfect detection, and provides an independent-validation estimator. The full corruption target still requires an exhaustive measured channel set, credible policy assignment, informative detection and amount restrictions, statistical uncertainty, and joint evidence on service quality, disruption and implementation resources.

\section*{References}
Benjamin A. Olken (2007), \emph{Monitoring Corruption: Evidence from a Field Experiment in Indonesia}, Journal of Political Economy 115(2), 200--249, \url{https://doi.org/10.1086/517935}. The source uses independent engineering assessments of project expenditure and motivates separating monitoring from measurement. The channel model and sharp detection bounds above are restricted derivations, not universal interpretations of that experiment.

\end{document}
